% from main.tex
%%%%%%%%%%  scan idx 269  |  folio 256  %%%%%%%%%%
plat. Grâce au critère de platitude par fibres (EGA IV 11.3.10), on en conclut que \(E^a/B^a\) est \(f\)-plat au voisinage de \(x\), cqfd.

\textbf{Corollaire 4.6.1.} Soient \(f:X\to S\) un morphisme localement de présentation finie et plat, \(x\in X\), \(s=f(x)\), \(E\in\Ob D^-(X)\), \([a,b]\) un intervalle de \(\mathbb Z\). Notons \(i_s:X_s=X\times_S\Spec(k(s))\to X\) la projection canonique, et posons \(E_{X_s}=\mathbf Li_s^*(E)\). On a alors
\[
\begin{gathered}
 \parfamp_f(E)\subset[a,b]\text{ au voisinage de }x\text{ et }\tamp(E_{X_s})_x\subset[a,b]\\
 \Longleftrightarrow \parfamp(E)\subset[a,b]\text{ au voisinage de }x.
\end{gathered}
\]
Par suite, \(E\) est parfait si et seulement si \(E\) est parfait rel. à \(f\) et \(E_{X_s}\) parfait pour tout \(s\in S\).

\emph{Preuve.} Appliquer (4.6) avec \(X=Y\), \(f=\Id_X\).

\textbf{Remarque 4.6.2.} De (4.6.1) et (I 5.10) résulte que, si \(f\) est localement de présentation finie, plat et à fibres des schémas réguliers, alors \(E\) parfait rel. à \(f\) équivaut à \(E\) parfait (cette équivalence avait déjà été obtenue (4.4) pour \(f\) lisse).

\textbf{Proposition 4.7.} Soient \(S\) un schéma, et pour \(i=1,2\), \(f_i:X_i\to Y_i\) un \(S\)-morphisme localement de type fini, \(E_i\in\Ob D(X_i)\), \([a_i,b_i]\) un intervalle de \(\mathbb Z\). Supposons \(X_1\) et \(X_2\), ainsi que \(Y_1\) et \(Y_2\) tor-indépendants sur \(S\) (1.5). Posons \(a=a_1+a_2\), \(b=b_1+b_2\), \(f=f_1\times_Sf_2\). Alors on a :
\begin{enumerate}[label=\alph*)]
\item Si \(\parfamp_{f_1}(E_1)\subset[a_1,b_1]\), et si \(E_2\) est pseudo-cohérent rel. à \(f_2\) et acyclique en degré extérieur à \([a_2,b_2]\), \(E_1\otimes_S^{\mathbf L}E_2\) (cf. (1.6)) est pseudo-cohérent rel. à \(f\) et acyclique en degré
